\documentclass[10pt]{article}
\usepackage[T1]{fontenc}
\usepackage[utf8]{inputenc}
\usepackage[french]{babel}
\usepackage{amsmath,amssymb,mathrsfs,mathpazo}
\usepackage{geometry}
\geometry{a4paper,margin=2cm}
\usepackage{graphicx,array,tabularx,tabularray,colortbl,multirow,multicol,booktabs,enumitem,pifont,ulem,fancybox,diagbox,cancel}
\usepackage{pstricks,pstricks-add,pst-plot,pst-tree,pst-node,pst-text,pst-func,pst-math,pst-bspline,pst-3dplot,pst-eucl,pst-solides3d}
\let\pscalculate\undefined
\usepackage{pst-bezier}
\usepackage{tikz,pgfplots,tkz-tab}
\usetikzlibrary{arrows,arrows.meta,patterns,calc,positioning,trees,decorations.markings,shapes.geometric,shapes.misc,shapes.arrows,intersections,decorations.pathmorphing,decorations.pathreplacing,backgrounds,shadows,babel}
\pgfplotsset{compat=1.18}
\usepgfplotslibrary{fillbetween}
\usepackage[np]{numprint}
\usepackage{listings,makecell,pas-tableur,twoopt}
\let\barre\undefined
\usepackage[pstricks]{tablvar}
\definecolor{darkgreen}{RGB}{0,100,0}
\providecommand{\red}{\color{red}}
\providecommand{\blue}{\color{blue}}
\providecommand{\green}{\color{green}}
\providecommand{\black}{\color{black}}
\providecommand{\white}{\color{white}}
\providecommand{\gray}{\color{gray}}
\providecommand{\cyan}{\color{cyan}}
\providecommand{\magenta}{\color{magenta}}
\providecommand{\yellow}{\color{yellow}}
\providecommand{\violet}{\color{violet}}
\providecommand{\orange}{\color{orange}}
\providecommand{\euro}{\texteuro}
\providecommand{\EUR}{\texteuro}
\providecommand{\ieme}{\textsuperscript{e}}
\providecommand{\degres}{\ensuremath{^\circ}}
\providecommand{\degree}{\ensuremath{^\circ}}
\providecommand{\vv}[1]{\overrightarrow{#1}}
\providecommand{\Coord}[1]{#1}
\providecommand{\decosix}{}
\providecommand{\starredbullet}{\ensuremath{\bullet}}
\pagestyle{empty}
\setlength{\parindent}{0pt}
\usepackage{adjustbox}\DeclareUnicodeCharacter{2236}{:}\DeclareUnicodeCharacter{2190}{\ensuremath{\leftarrow}}
\providecommand{\R}{\mathbb{R}}
\providecommand{\N}{\mathbb{N}}
\providecommand{\Z}{\mathbb{Z}}
\providecommand{\Q}{\mathbb{Q}}
\providecommand{\D}{\mathbb{D}}
\providecommand{\C}{\mathbb{C}}
\providecommand{\e}{{\mathrm{e}}}
\providecommand{\vect}[1]{\overrightarrow{#1}}
\providecommand{\vectt}[1]{\overrightarrow{\mathrm{#1}}}
\providecommand{\vvt}[1]{\overrightarrow{\mathrm{#1}}}
\providecommand{\np}[2][]{{#2~#1}}
\providecommand{\vv}[1]{{\overrightarrow{#1}}}
\renewcommand{\d}{{\mathrm{d}}}
\providecommand{\degres}{{^\circ}}
\providecommand{\degre}{{^\circ}}
\providecommand{\blue}{\color{blue}}
\providecommand{\red}{\color{red}}
\providecommand{\enskip}{\quad}
\providecommand{\llbracket}{\lbrack\!\lbrack}
\providecommand{\rrbracket}{\rbrack\!\rbrack}
\providecommand{\euro}{{\text{€}}}
\providecommand{\no}{{\text{n°}}}
\providecommand{\interval}[2]{{[#1;#2]}}
\providecommand{\intervFO}[2]{{[#1;#2[}}
\providecommand{\dotfill}{\ldots}

\providecommand{\ds}{\displaystyle}
\providecommand{\textding}[1]{#1}
\providecommand{\Oij}{\(\mathrm{O};\vec{\imath},\vec{\jmath}\)}
\providecommand{\Oijk}{\(\mathrm{O};\vec{\imath},\vec{\jmath},\vec{k}\)}
\providecommand{\Ouv}{\(\mathrm{O};\vec{u},\vec{v}\)}
\providecommand{\intervFF}[2]{\ensuremath{[#1;#2]}}
\renewcommand{\euro}{\ifmmode\text{\texteuro}\else\texteuro\fi}
\DeclareUnicodeCharacter{00B0}{\ifmmode^{\circ}\else\textdegree\fi}
\DeclareUnicodeCharacter{20AC}{\ifmmode\text{\texteuro}\else\texteuro\fi}
% Énoncés : Éducation nationale. Transcription et corrigés : APMEP, Denis Vergès et contributeurs. Découpage et classement : CoopMaths, Sébastien Lozano, Éric Elter et contributeurs. Adaptation de la présentation : Maths Explorer.
% CC BY-SA : https://coopmaths.fr/www/about/
% https://coopmaths.fr/alea/static/bac/2024/tex/bac_2024_06_sujet2_metropole_devoile_3.tex
\begin{document}
\noindent\textit{APMEP / CoopMaths. CC BY-SA. Présentation adaptée par Maths Explorer.}\par\medskip
\medskip

Soit $a$ un nombre réel strictement supérieur à 1.

On considère la suite $\left(u_n\right)$ définie par $u_0 = a$ et, pour tout entier naturel $n$ : \[u_{n+1} = u_n^2 - 2u_n + 2.\]

On admet que pour tout entier naturel $n$,\: $u_n > 1$.

L'objectif de cet exercice est d'étudier la suite $\left(u_n\right)$ pour différentes valeurs du nombre réel $a$.

\bigskip

\textbf{Partie A : étude de la suite \boldmath $\left(u_n\right)$ \unboldmath dans le cas }\boldmath$1 < a < 2$\unboldmath 

\medskip

\begin{enumerate}
\item 
	\begin{enumerate}
		\item Montrer que, pour tout entier naturel $n$, on a : $u_{n+1} - 2 = u_n\left(u_n - 2\right)$.
		\item Montrer que, pour tout entier naturel $n$, on a : $u_{n+1} - u_n = \left(u_n - 1\right)\left(u_n - 2\right)$.
	\end{enumerate}
\item Dans cette question, on pourra utiliser les égalités établies dans la question précédente.
	\begin{enumerate}
		\item En utilisant un raisonnement par récurrence démontrer que, pour tout entier naturel $n : u_n < 2$.
		\item Montrer que la suite $\left(u_n\right)$ est convergente et déterminer sa limite. 
	\end{enumerate}
\end{enumerate}
	
	
\textbf{Partie B : étude dans le cas particulier }\boldmath $a = 2$\unboldmath

\medskip

\begin{enumerate}
\item ~\\[-25pt]

\begin{minipage}{0.5\linewidth}
On donne ci-contre la fonction \texttt{u} écrite en langage Python.

Déterminer les valeurs renvoyées par le programme lorsque l'on saisit \texttt{u(2,1)} et \texttt{u(2,2)} dans la console Python.
\end{minipage}
\hfill
\begin{minipage}{0.45\linewidth}
{\ttfamily
\begin{tabular}{|l|}\hline
\pmb{def} u(a,n) : \\
\qquad u=a\\
\qquad \pmb{for} k \pmb{in} range(n) :\\
\qquad \quad  u=u**2-2*u+2\\
\qquad \pmb{return} u\\ \hline
\end{tabular}
}
\end{minipage}

\item Quelle conjecture peut-on formuler concernant la suite $\left(u_n\right)$ dans le cas où $a = 2$ ? 

On admettra ce résultat sans démonstration.
\end{enumerate}

\bigskip

\textbf{Partie C : étude dans le cas général}

\medskip

\begin{enumerate}
\item On considère la suite $\left(v_n\right)$ définie, pour tout entier naturel $n$, par $v_n = \ln \left(u_n - 1\right)$.
	\begin{enumerate}
		\item Montrer que la suite $\left(v_n\right)$ est une suite géométrique de raison 2 dont on précisera le premier terme en fonction de $a$.
		\item En déduire que, pour tout entier naturel $n,\: u_n = 1 + \e^{2^n \times \ln (a-1)}$.
	\end{enumerate}
\item Déterminer, suivant les valeurs du réel $a$ strictement supérieur à 1, la limite de la suite $\left(u_n\right)$.
\end{enumerate}

\bigskip

\end{document}
